% 题干转写：完成
% 解析转写：完成
% 图形核对：题面与答案均无独立图件；补充解析辅助图
\section{第39届复赛第5题：快慢拉动双弹簧断裂模型}

\subsection{题干}
两根相同的理想轻弹簧，劲度系数均为 $k$、自然长度均为零；中间以质量为 $m$ 的小球相连。弹簧1上端固定在天花板，系统自然下垂且初始静止。任一弹簧拉伸到临界长度（大于 $mg/k$）即断裂。对弹簧2下端施加随时间变化的外力 $F(t)$。

设 $F(t)=0$（$t<0$），$F(t)=\alpha t$（$t\ge0$），其中 $\alpha>0$ 表示加载快慢。求：（1）弹簧1长度的运动方程；（2）$x_1(t),x_2(t)$；（3）两根弹簧先断的条件及其与 $\alpha$ 的关系；（4）临界长度为 $L$ 时确保弹簧2先断的 $\alpha$ 条件，以及同时断裂的条件。

\begin{figure}[htbp]\centering
  \begin{tikzpicture}[>=Stealth,scale=0.98]
  \draw[very thick] (-1.5,2.5) -- (1.5,2.5);
  \foreach \x in {-1.35,-1.15,...,1.35} {\draw[thin] (\x,2.5) -- ++(0.15,0.11);}
  \draw[thick] (0,2.5) -- (0,2.18);
  \draw[thick,decorate,decoration={zigzag,segment length=5pt,amplitude=2.2pt}] (0,2.18) -- (0,0.76);
  \node[left] at (-1.82,1.5) {弹簧1};
  \shade[ball color=gray!35] (0,0.55) circle (0.21);
  \node[right] at (0.24,0.55) {$m$};
  \draw[thick,decorate,decoration={zigzag,segment length=5pt,amplitude=2.2pt}] (0,0.34) -- (0,-1.06);
  \node[right] at (1.72,-0.35) {弹簧2};
  \draw[thick] (0,-1.06) -- (0,-1.35);
  \draw[->,very thick] (0,-1.35) -- (0,-2.1) node[right] {$F(t)=\alpha t\quad(t\geq0)$};
  \draw[->,thick] (0,0.42) -- (0,-0.35) node[right] {$mg$};
  \draw[<->,thin] (-1.1,2.48) -- (-1.1,0.75) node[midway,right] {$x_1(t)$};
  \draw[<->,thin] (0.84,0.34) -- (0.84,-1.06) node[midway,right] {$x_2(t)$};
  \node[below=0.3cm] at (0,-2.1) {解析辅助示意：零自然长度双弹簧的加载模型};
\end{tikzpicture}

  \caption{第39届复赛第5题的解析辅助加载模型。题面和答案均无独立图件，本图仅定义 $x_1,x_2$ 与外力方向。}\label{fig:cpho39q05model}
\end{figure}

\subsection{参考解答整理}
\paragraph{（1）运动方程。}
以弹簧1的长度 $x_1(t)$ 描述小球的竖直位置。牛顿第二定律与胡克定律给出
\[
\boxed{m\ddot{x}_1=-kx_1+mg+F(t).}
\]

\paragraph{（2）线性加载下的长度。}
令 $\omega^2=k/m$，以 $mg/k$ 平移坐标并代入 $F(t)=\alpha t$。由初始条件 $x_1(0)=mg/k$、$\dot{x}_1(0)=0$，得到
\[
\boxed{x_1(t)=\frac{mg}{k}+\frac{\alpha}{k}t-\frac{\alpha}{\omega k}\sin\omega t,}
\qquad
\boxed{x_2(t)=\frac{\alpha}{k}t.}
\]
后式直接由 $F(t)=kx_2(t)$ 得到。

\paragraph{（3）先断条件与快慢加载。}
两弹簧长度差为
\[
\Delta x=x_1-x_2=\frac{mg}{k}-\frac{\alpha}{\omega k}\sin\omega t.
\]
因此当 $\alpha<mg\omega=g\sqrt{km}$ 时，任何时刻均有 $\Delta x>0$，弹簧1必先断。若 $\alpha\ge mg\omega$，弹簧1先断需满足
\[
\sin(\omega t_0)<\frac{mg\omega}{\alpha};
\]
弹簧2先断则需满足相反的不等式，且必须有 $\alpha>mg\omega$。图\ref{fig:cpho39q05phase}将该长度差的相位判定画出：加载越快，满足弹簧2先断的相位区间越宽，这解释了快速拉动下弹簧2更可能断裂的现象。
\begin{figure}[htbp]\centering
  \begin{tikzpicture}[>=Stealth,x=1.15cm,y=1.05cm]
  \draw[->] (-0.2,0) -- (6.7,0) node[right] {$\phi=\omega t$};
  \draw[->] (0,-1.55) -- (0,1.75) node[above] {$\Delta x=x_1-x_2$};
  \draw[thick,blue!65!black,domain=0:6.28,samples=160] plot (\x,{0.45-0.95*sin(\x r)});
  \draw[densely dashed] (0,0.45) -- (6.35,0.45) node[right] {$mg/k$};
  \draw[densely dashed] (0,0) -- (6.35,0);
  \foreach \x/\lab in {1.57/{\pi/2},3.14/{\pi},4.71/{3\pi/2},6.28/{2\pi}} {
    \draw[densely dotted,gray] (\x,-1.45) -- (\x,1.55);
    \node[below] at (\x,0) {$\lab$};
  }
  \fill[blue!65!black] (0.49,0) circle (1.5pt);
  \fill[blue!65!black] (2.65,0) circle (1.5pt);
  \node[below left,align=center] at (0.49,-0.05) {$\Delta x=0$};
  \node[above right,align=center] at (2.65,0.04) {$\Delta x=0$};
  \draw[<->,thick] (0.52,-1.24) -- (2.62,-1.24) node[midway,below] {弹簧2可能先断的相位区间};
  \node[blue!65!black,above] at (4.55,0.95) {$\Delta x=\frac{mg}{k}-\frac{\alpha}{\omega k}\sin\phi$};
  \node[below=0.65cm] at (3.15,-1.5) {解析辅助示意：$\alpha>mg\omega$ 时的长度差与相位};
\end{tikzpicture}

  \caption{第39届复赛第5题的解析辅助相位图：长度差符号决定先断弹簧。}\label{fig:cpho39q05phase}
\end{figure}

\paragraph{（4）给定临界长度 $L$。}
弹簧2达到临界长度的时刻为 $t_0=kL/\alpha$。要确保其先断，应满足
\[
\boxed{\sin\left(\frac{kL}{\alpha}\sqrt{\frac{k}{m}}\right)
>\frac{mg}{\alpha}\sqrt{\frac{k}{m}}.}
\]
存在同时断裂的情形，其临界条件只需将上式的不等号改为等号。

\subsection{转写依据}
题干依据 \texttt{CPHO39-Q05.md} 第3页，答案依据第20至23页与 \texttt{CPHO39-Q05\_formula\_check\_v1.csv}。题面和答案均无独立图件；图\ref{fig:cpho39q05model}、图\ref{fig:cpho39q05phase}均为对应推导关系的自绘辅助图。本题四问分值依次为3、19、13、5分。
